\writeSectionFrame{\presentationTitle}{Singleton}
\begin{frame}{Steckbrief}
	\begin{itemize}
 		\item Deutscher Name: Singleton
 		\item Alternativer deutscher Name: Einzelstück (eher selten benutzt)
 		\item Englischer Name: Singleton
 		\item Gruppe: Erzeugungsmuster
 	\end{itemize}
\end{frame}

\begin{frame}{Singleton}
	\begin{block}{Zweck}
 		Sichere, dass eine Klasse genau ein Exemplar besitzt, und stelle einen globalen Zugriffspunkt darauf zur Verfügung.
 	\end{block}
\end{frame}

\frameWithImageOnly{\BILDERPFAD/singleton1.jpg}
\note{
	Das erste einfache Muster ist der Singleton. Dieses Muster wird eingesetzt,
	wenn es von einer Klasse nur ein einziges Objekt geben darf. Auf dieses Objekt kann dann vom gesamten
	Programm aus zugegriffen werden.Wir erinnern uns an Spring: Dort wurden standardmäßig 
	alle von Spring verwalteten Objekte als Singletons erstellt und verwaltet. Als
	Singleton verwalten werden können z.B. Fassaden, Dao-Klassen oder Klassen mit der 
	Geschäftslogik.  
}

\frameWithImageOnly{\BILDERPFAD/singleton2.jpg}
\note{
	 Beim Zustand der Singleton-Objekte ist Vorsicht geboten, denn sie dürfen natürlich keinen
	 objektspezifischen Zustand besitzen. 
}

\begin{frame}{Klassendiagramm}
	\begin{center}
 		\includegraphics[scale=0.65]{\BILDERPFAD/singleton3.png}
 	\end{center}
\end{frame}

\note{
	Das Klassendiagramm ist das einfachste der Entwurfsmuster. Der Singleton verwaltet das einzige
	existierende Objekt selbst und somit können nicht mehr mehrere Objekte erzeugt werden.
}

\begin{frame}[fragile]{Einzige Instanz}
	\begin{exampleblock}{ConfigurationSingleton.java}
		\begin{lstlisting}
public final class ConfigurationSingleton {
	private static final String CONFIGURATION_FILE = 
        "configuration.properties";
	private Map<String, String> configurationMap;
	private static ConfigurationSingleton configuration;
			    
		\end{lstlisting}
 	\end{exampleblock}
\end{frame}

\note{
	Zuerst definieren wir in der Singleton-Klasse die einzige Instanz als statisches Attribut
}

\begin{frame}[fragile]{Methode getInstance()}
	\begin{exampleblock}{ConfigurationSingleton.java}
		\begin{lstlisting}
public static ConfigurationSingleton getInstance() {
    if (configuration == null) {
        configuration = new ConfigurationSingleton();
    }
    return configuration;
}		    
		\end{lstlisting}
 	\end{exampleblock}
\end{frame}

\note{
	Dann brauchen wir eine Methode, die das einzige Objekt zurückgibt, so dass man damit arbeiten kann.
	Wir haben aber ein Problem: Wir stellen nicht sicher, dass es tatsächlich nur eine einzige Instanz
	gibt. Aber wie erzeugen wir diese Instanz? Das tun wir in genau dieser Methode getInstance(). Wir sehen nach, ob das Objekt bereits existiert und
	wenn nicht, wird es instantiiert.
}

\begin{frame}[fragile]{Verwendung}
	\begin{exampleblock}{ConfigurationSingleton.java}
		\begin{lstlisting}
private static void singletonExample() {
    ConfigurationSingleton configuration1 = 
        ConfigurationSingleton.getInstance();
 
    configuration1.addValue("background", "black");
    configuration1.addValue("margin", "10mm");
    
    ConfigurationSingleton configuration2 = 
        ConfigurationSingleton.getInstance();
    
    configuration2.addValue("italic", "true");
    configuration2.addValue("bold", "false");
    
    configuration1.writeToFile();
    configuration2.writeToFile();
}
	 	\end{lstlisting}
 	\end{exampleblock}
\end{frame}

\note{
	Wir haben jetzt tatsächlich nur ein einziges Objekt, wenn wir mit der Methode getInstance() arbeiten. Aber
	wir erzwingen noch nicht, dass tatsächlich nur ein einziges Objekt erzeugt wird.
}

\begin{frame}[fragile]{Singleton}
	\begin{exampleblock}{ConfigurationSingleton.java}
		\begin{lstlisting}
private ConfigurationSingleton() {
    initializeFile();
    configurationMap = new HashMap<String, String>();
    loadConfigurationFromFile();
}
	 	\end{lstlisting}
 	\end{exampleblock}
\end{frame}

\note{
    Indem wir den Konstruktor privat setzen, verhindern wir jede weitere Objekterzeugung und somit können keine weiteren Objekte mehr erzeugt werden.
}

